\subsection{Edge-Preserving Filtering}

Filtering is an operation to reduce the impact of sensor noise from limited resolution and other random effects on measured quantities.
These deviations propagate through the feature image conversion.
The feature images, introduced in Section~\ref{sec:feature_images}, build upon information information about changes in the geometry of sensed objects.
Therefore, sharp changes, e.g.~edges and corners, shall be preserved by the filter, requiring edge-preserving filters.
Each tested filter uses OpenCV's\cite{opencv_library} proven and readily available implementation.

\subsubsection{Median Blur}

Median Blur, introduced to image processing by Frieden\cite{frieden_new76}, is a well established technique to reduce salt-and-pepper noise and is an effective edge-preserving filter.
The filter modifies each pixel of the image equally.
A window of $n \times m$ pixels with $n,m \in \mathbb{N}_{2k + 1}$ is centered at the pixel.
The median of all pixels in that window is calculated and stored as the new value.
Floating point operations are not required and the filter is implementable with $\mathcal{O}(n)$\cite{huang_ieee79} time complexity.
$n$ and $m$ are the controllable parameters.
OpenCV's\cite{opencv_library} implementation additionally utilizes \acrshort{SIMD} (\acrlong{SIMD}) instructions for increased processing speed.

\subsubsection{Bilateral Filter}

Bilateral filtering, introduced by Tomasi and Manduchi\cite{tomasi_iccv98}, considers two factors in the filtering process.
\emph{Spatial closeness}, similar to median filtering, and \emph{similarity} of the values are combined in the bilateral filter.
Its application is convolutional, similar to the median blur and other filters.
Each neighbouring pixel's effect on the central pixel is weighted by a geometric \emph{closeness} and \emph{similarity} function.
Similarity can be based on a classical distance norm of photometric values, like brightness, or be more sophisticated, like perceived similarity.
A common function for both aspects is a Gaussian kernel with the Euclidean norm as distance measure.
Figure~\ref{fig:bilateral_filter} provides an example application of the bilateral filter on a single channel, two-dimensional signal with the Gaussian function for \emph{closeness} and \emph{similarity}.
\begin{figure}[b!]
    \begin{subfigure}[b]{0.3\linewidth}
        \includegraphics[width=\linewidth]{chapter04/img/bilateral1.png}
        \caption{}\label{fig:bilateral_1}
    \end{subfigure}
    \begin{subfigure}[b]{0.3\linewidth}
        \includegraphics[width=\linewidth]{chapter04/img/bilateral2.png}
        \caption{}\label{fig:bilateral_2}
    \end{subfigure}
    \begin{subfigure}[b]{0.3\linewidth}
        \includegraphics[width=\linewidth]{chapter04/img/bilateral3.png}
        \caption{}\label{fig:bilateral_3}
    \end{subfigure}
    \caption[Bilateral filtering visualized]{\emph{Bilateral filtering visualized.} The figures show how the bilateral filter works on a single channel step function. Figure~\ref{fig:bilateral_1} shows the original signal with random noise added. The weighting of its neighbouring pixels is computed for a pixel in the center. The weights are visualized in Figure~\ref{fig:bilateral_2}. The lower values of the step are neglected due to their lack of similarity regardless of closeness. Figure~\ref{fig:bilateral_3} shows the result of the full convolution of the filter, yielding a smoothed signal without blurred edge. Images adopted from~\cite{tomasi_iccv98}.}\label{fig:bilateral_filter}
\end{figure}
The filter is dependant on the functions for \emph{similarity} and \emph{closeness} and the available parameters are defined through these functions.
The Gaussian case is controlled by the parameters $\sigma_d$, the standard deviation of \emph{closeness} and $\sigma_r$, the standard deviation of \emph{similarity}.
The mathematical formulation --- based on the original publication\cite{tomasi_iccv98} --- for the \emph{closeness}~$c$ in the Gaussian case of the bilateral filter is:
\begin{equation}
\begin{aligned}
    c(\xi, \vec{x})&= \exp\left(-\frac{1}{2}{\left(\frac{d(\xi, \vec{x})}{\sigma_d} \right)}^2 \right) \\
    d(\xi, \vec{x})&= d(\xi - \vec{x}) = \lnorm{\xi - \vec{x}} \text{.}
\end{aligned}
\end{equation}
This weighs any point's $\xi$ relevance from the input signal to the filtered point $\vec{x}$ based on it's distance.
The definition of the \emph{similarity}~$s$ follows the same principle:
\begin{equation}
\begin{aligned}
    s(\xi, \vec{x})&= \exp\left(-\frac{1}{2}{\left(\frac{\theta(\mathbf{f}(\xi), \mathbf{f}(\vec{x}))}{\sigma_r}\right)}^2\right) \\
    \theta(\phi, \mathbf{f})&= \theta(\phi - \mathbf{f}) = \lnorm{\phi - \mathbf{f}}\text{.}
\end{aligned}
\end{equation}
The function $\mathbf{f}(\vec{x})$ is the value of the signal at the position $\vec{x}$ and $\theta(\phi, \mathbf{f})$ weighs these values based on their Euclidean distance.
To compute the filtered value on any point $\vec{x}$, these functions need to be applied to the whole signal.
The continuous formulation is given here, but the filter is a discrete convolution for images.
\begin{equation}
\begin{aligned}
    h(\vec{x}) &= \frac{1}{k(\vec{x})} \int_{-\infty}^{\infty} \int_{-\infty}^{\infty} \mathbf{f}(\xi) c(\xi, \vec{x}) s(\mathbf{f}(\xi), \mathbf{f}(\vec{x})) d\xi \\
    \text{with the normalization:} \\
    k(\vec{x}) &= \int_{-\infty}^{\infty} \int_{-\infty}^{\infty} c(\xi, \vec{x}) s(\mathbf{f}(\xi), \mathbf{f}(\vec{x})) d\xi
\end{aligned}
\end{equation}
Higher values of $\sigma_d$ and $\sigma_r$ result in stronger filtering of the result.
The bilateral filter is a single pass filter, but computationally expensive and slower than median blur.
